%!TEX root = ../main.tex
% RESUMEN EN ESPAÑOL — MARKETLENS
% Revisar nuevamente al cerrar encuesta, entrevista y apéndices.
% Las palabras clave se configuran en main.tex con \palabrasclave{...}.

\begin{resumen}

Las plataformas de compraventa digital facilitan el intercambio de productos entre usuarios, pero también exponen a los compradores a publicaciones engañosas, información incompleta y dificultades para valorar la confiabilidad de los vendedores. Aunque los marketplaces disponen de mecanismos propios de reputación y protección, las señales relevantes pueden encontrarse dispersas o resultar insuficientes para interpretar el riesgo de una operación.

Este trabajo presenta MarketLens, un prototipo de plataforma de análisis preventivo de publicaciones y vendedores. Mediante una metodología de investigación aplicada y desarrollo incremental, se implementaron una interfaz web, una extensión de navegador, una interfaz de programación de aplicaciones desarrollada con Java y Spring Boot y una base de datos PostgreSQL. El sistema reúne atributos observables, aplica reglas ponderadas, compara precios cuando existen referencias suficientemente comparables y comunica un índice de riesgo acompañado de factores explicativos y una medida de cobertura de la información.

Las pruebas funcionales permitieron comprobar la integración de componentes sobre publicaciones seleccionadas de Mercado Libre. En uno de los casos documentados, la extensión mostró un índice de 40 sobre 100 con cobertura parcial de dos de tres grupos de información. Asimismo, se registró una ejecución de 62 pruebas automatizadas sin fallos ni errores. Las pruebas de comparación de precios evidenciaron la necesidad de evitar conclusiones cuando las referencias disponibles no son suficientes.

MarketLens aporta una base técnica para reunir y explicar indicadores preventivos. Sin embargo, los resultados acreditan el funcionamiento del prototipo en los casos ensayados, no su capacidad estadísticamente validada para identificar fraudes reales. La incorporación de reportes colaborativos y una validación más amplia permanecen como líneas de trabajo futuro.

\end{resumen}
